% Generated draft: requires mathematical and source review.
\documentclass[11pt]{article}
\usepackage[margin=25mm]{geometry}
\usepackage{fontspec}
\setmainfont{Latin Modern Roman}
\usepackage{amsmath,amssymb,mathtools,unicode-math}
\setmathfont{Latin Modern Math}
\usepackage{enumitem,needspace,xurl,hyperref}
\setlength{\parindent}{0pt}
\setlength{\parskip}{5pt}
\setlength{\emergencystretch}{3em}
\newcommand{\N}{\mathbb{N}}
\newcommand{\Z}{\mathbb{Z}}
\newcommand{\Q}{\mathbb{Q}}
\newcommand{\R}{\mathbb{R}}
\newcommand{\C}{\mathbb{C}}
\newcommand{\F}{\mathbb{F}}
\newcommand{\A}{\mathbb{A}}
\newcommand{\PP}{\mathbb{P}}
\newcommand{\E}{\mathbb{E}}
\newcommand{\eps}{\varepsilon}
\newcommand{\dd}{\,\mathrm{d}}
\begin{document}
\section*{UnsolvedMath ID: 30004135}
\textbf{Source problem code:} catalogue entry
\subsection{Definitions and mathematical statement}
An edge ordering of a finite graph \(G=(V,E)\) is a strict total order \(\prec\) on \(E\). Write \(n=|V|\) and \(m=e(G)=|E|\). An increasing path of length \(\ell\) consists of distinct vertices \(v_0,\ldots,v_\ell\) such that \(v_{i-1}v_i\in E\) for \(1\le i\le\ell\) and
\[
 v_0v_1\prec v_1v_2\prec\cdots\prec v_{\ell-1}v_\ell.
\]
Thus length counts edges; a single vertex gives a path of length zero.

The supplied record does not explicitly specify whether graphs are finite, simple, and undirected. Subject to confirming this conventional interpretation, and restricting to \(n\ge1\) so that the density is defined, the question is: do there exist absolute constants \(\varepsilon>0\) and \(c>0\) such that, for every finite simple undirected graph \(G\) with at least one vertex and every strict total ordering of its edges, there is an increasing path of length \(\ell\) satisfying
\[
 \ell\ge c\left(\frac{m}{n}\right)^{1/2+\varepsilon}?
\]
The constants must be independent of the graph and its edge ordering. This makes explicit the uniform interpretation of the source's \(\Omega\)-notation. The parameter \(m/n\) is half the average degree. If \(m=0\), the bound is zero and a single vertex suffices. The graph with no vertices is outside this formulation because \(m/n\) is undefined.

The graph conventions and the edge-count convention for path length require confirmation against the source. The supplied registry labels the problem open and its statement \texttt{corrected\_verified}; these are supplied curation labels, not independent mathematical verification.
\subsection{Short English statement}
Is there a fixed positive improvement over the square-root exponent such that every edge-ordered graph has an increasing path whose length is at least a universal constant times \((e(G)/|V(G)|)^{1/2+\varepsilon}\)? The explicit formulation uses finite simple undirected graphs with at least one vertex, pending confirmation of these unstated source conventions.
\subsection{Sources}
\begin{itemize}
\item Local UnsolvedMath registry, supplied complete record for permanent ID \(30004135\), titled ``Long Increasing Paths in Edge-Ordered Graphs,'' source problem code \texttt{catalogue entry}, in the Oberwolfach Reports Open Problems collection. The record gives proposed year 2019 and names no proposer. Its statement-verification note says that mathematical content was preserved while notation, typography, and sentence structure were repaired for publication; this claim is supplied metadata.
\item Supplied report citation: \emph{Combinatorics, Probability and Computing} (2020), Oberwolfach Reports, pp.~1113--1189, DOI \texttt{10.4171/owr/2019/19}. \url{https://doi.org/10.4171/owr/2019/19}. The background describes the collection as drawn from reports documenting mathematical workshops and their open-problem discussions. No page locating this particular problem within the report is supplied.
\item Supplied literature reference, attributed to Bucic--Kwan--Pokrovskiy--Sudakov--Tran--Wagner: \url{https://arxiv.org/abs/1809.01468}.
\item Supplied literature reference labeled ``Published paper'': \url{https://people.math.ethz.ch/~sudakovb/monotone-path-complete-graph.pdf}.
\item Literature assessment dated exactly \texttt{2026-08-22}: the supplied curation labels the problem open and states that nearly linear monotone paths are proved for complete graphs and reasonably dense edge-ordered graphs, while the fixed-average-degree sparse regime in the stated problem remains unresolved. This is supplied dated curation derived from the OWR material, not an independently verified assessment. Its time-sensitive status and cited claims require independent verification.
\end{itemize}
\end{document}
